\section{Exponentials and logarithms}

\begin{convention}
\label{conv.fps.exp.K-Q-alg}Throughout this section (unless otherwise stated), we assume
that $K$ is not just a commutative ring, but actually a commutative
$\mathbb{Q}$-algebra.
\end{convention}

\subsection{Definitions}

\begin{definition}
\label{def.fps.exp-log}
\lean{PowerSeries.logbar, PowerSeries.expbar}
\leantarget
\leanok
Define three FPS $\exp$, $\overline{\log}$ and
$\overline{\exp}$ in $K\left[\left[x\right]\right]$ by
\begin{align*}
\exp &  :=\sum_{n\in\mathbb{N}}\dfrac{1}{n!}x^{n},\\
\overline{\log}  &  :=\sum_{n\geq1}\dfrac{\left(-1\right)^{n-1}}{n}x^{n},\\
\overline{\exp}  &  :=\exp-1=\sum_{n\geq1}\dfrac{1}{n!}x^{n}.
\end{align*}

(The last equality sign here follows from $\exp=\sum_{n\in\mathbb{N}}\dfrac
{1}{n!}x^{n}=\underbrace{\dfrac{1}{0!}}_{=1}\underbrace{x^{0}}_{=1}
+\sum_{n\geq1}\dfrac{1}{n!}x^{n}=1+\sum_{n\geq1}\dfrac{1}{n!}x^{n}$.)
\end{definition}

\subsection{Helpers for derivative computations}

\begin{lemma}[The series $(1+x)^{-1}$]
\label{lem.fps.invOnePlusX}
\lean{PowerSeries.invOnePlusX}
\leanhelper
\leanok
Define the FPS
\[
\iota := \sum_{n \in \mathbb{N}} (-1)^n x^n.
\]
This is the geometric series for $(1+x)^{-1}$; indeed, $\iota \cdot (1+x) = 1$.
\end{lemma}

\begin{proof}
Direct verification: the product is computed coefficient-wise and shown to equal~$1$.
\end{proof}

\begin{lemma}[Derivative of $\overline{\log}$]
\label{lem.fps.derivative-logbar}
\lean{PowerSeries.derivative_logbar}
\leanhelper
\leanok
We have $\overline{\log}' = \iota = (1+x)^{-1}$.
\end{lemma}

\begin{proof}
\leanok
By computing the derivative of $\overline{\log} = \sum_{n \ge 1} \frac{(-1)^{n-1}}{n} x^n$
term-by-term:
$\overline{\log}' = \sum_{n \ge 1} (-1)^{n-1} x^{n-1} = \sum_{n \ge 0} (-1)^n x^n = \iota$.
\end{proof}

\begin{lemma}[Derivative of $\overline{\exp}$]
\label{lem.fps.derivative-expbar}
\lean{PowerSeries.derivative_expbar}
\leanhelper
\leanok
We have $\overline{\exp}' = \exp$.
\end{lemma}

\begin{proof}
\leanok
Since $\overline{\exp} = \exp - 1$, we have $\overline{\exp}' = \exp' - 0 = \exp' = \exp$,
where the last equality is the standard fact $\exp' = \exp$.
\end{proof}

\begin{lemma}[Substitution of $(1+x)^{-1}$]
\label{lem.fps.invOnePlusX-subst}
\lean{PowerSeries.invOnePlusX_subst_eq_inv}
\leanhelper
\leanok
Let $g \in K\llbracket x \rrbracket$ with $[x^0]g = 0$.
Then $\iota \circ g = (1+g)^{-1}$.
\end{lemma}

\begin{proof}
\leanok
We have $(\iota \circ g) \cdot (1 + g)
= (\iota \cdot (1+x)) \circ g = 1 \circ g = 1$.
Since $[x^0](1+g) \ne 0$, this gives $\iota \circ g = (1+g)^{-1}$.
\end{proof}

\subsection{The exponential and the logarithm are inverse}

\begin{proposition}
\label{prop.fps.exp-log-der}
\lean{PowerSeries.derivative_exp_comp, PowerSeries.derivative_expbar_comp, PowerSeries.derivative_logbar_comp}
\leantarget
\leanok
Let $g\in K\left[\left[x\right]\right]$
with $\left[x^{0}\right]g=0$. Then: \medskip

\textbf{(a)} We have
\[
\left(\overline{\exp}\circ g\right)^{\prime}=\left(\exp\circ g\right)
^{\prime}=\left(\exp\circ g\right)\cdot g^{\prime}.
\]

\textbf{(b)} We have
\[
\left(\overline{\log}\circ g\right)^{\prime}=\left(1+g\right)
^{-1}\cdot g^{\prime}.
\]
\end{proposition}

\begin{proof}
\textbf{(a)} Let us first
show that $\overline{\exp}^{\prime}=\exp^{\prime}=\exp$. Indeed,
$\overline{\exp}=\exp-1$, so that $\exp=\overline{\exp}+1$ and therefore
\begin{align}
\exp^{\prime} &  =\left(\overline{\exp}+1\right)^{\prime}=\overline{\exp
}^{\prime}+\underbrace{1^{\prime}}_{=0}
=\overline{\exp}^{\prime}. \label{pf.prop.fps.exp-log-der.a.1}
\end{align}
Next, we recall that $\exp=\sum_{n\in\mathbb{N}}\dfrac{1}{n!}x^{n}$. Hence,
the definition of a derivative yields
\begin{align}
\exp^{\prime} &  =\sum_{n\geq1}\underbrace{n\cdot\dfrac{1}{n!}}%
_{\substack{=\dfrac{1}{\left(n-1\right)!}\\\text{(since }n!=n\cdot\left(
n-1\right)!\text{)}}}x^{n-1}=\sum_{n\geq1}\dfrac{1}{\left(n-1\right)
!}x^{n-1}=\sum_{n\in\mathbb{N}}\dfrac{1}{n!}x^{n}\nonumber\\
& \qquad\qquad\left(\text{here, we have
substituted }n\text{ for }n-1\text{ in the sum}\right) \nonumber\\
&  =\exp. \label{pf.prop.fps.exp-log-der.a.2}
\end{align}
Comparing this with (\ref{pf.prop.fps.exp-log-der.a.1}), we find
\begin{equation}
\overline{\exp}^{\prime}=\exp. \label{pf.prop.fps.exp-log-der.a.3}
\end{equation}

Now, we can apply the chain rule to $f=\overline{\exp}$
(since $\left[x^{0}\right]g=0$), and
thus obtain
\[
\left(\overline{\exp}\circ g\right)^{\prime}=\left(\underbrace{\overline
{\exp}^{\prime}}_{=\exp}\circ\,g\right)\cdot g^{\prime}=\left(\exp\circ
g\right)\cdot g^{\prime}.
\]
The same computation (but with $\overline{\exp}$ replaced by $\exp$) yields
$\left(\exp\circ g\right)^{\prime}=\left(\exp\circ g\right)\cdot
g^{\prime}$. Combining these two formulas, we obtain $\left(\overline{\exp
}\circ g\right)^{\prime}=\left(\exp\circ g\right)^{\prime}=\left(
\exp\circ g\right)\cdot g^{\prime}$. This proves part \textbf{(a)}.

\textbf{(b)} We have $\overline{\log}=\sum_{n\geq1}\dfrac{\left(-1\right)
^{n-1}}{n}x^{n}$. Thus,
\begin{align*}
\overline{\log}^{\prime} &  =\left(\sum_{n\geq1}\dfrac{\left(-1\right)
^{n-1}}{n}x^{n}\right)^{\prime}
=\sum_{n\geq1}\underbrace{\dfrac{\left(-1\right)^{n-1}}{n}n}_{=\left(
-1\right)^{n-1}}\underbrace{x^{\prime}}_{=1}x^{n-1}=\sum_{n\geq1}\left(
-1\right)^{n-1}x^{n-1}=\sum_{n\in\mathbb{N}}\left(-1\right)^{n}x^{n}
\end{align*}
(here, we have substituted $n$ for $n-1$ in the sum). On the other hand,
$\left(1+x\right)^{-1}
=\sum_{n\in\mathbb{N}}\left(-1\right)^{n}x^{n}$. Comparing these two
equalities, we find
\begin{equation}
\overline{\log}^{\prime}=\left(1+x\right)^{-1}.
\label{pf.prop.fps.exp-log-der.b.2}
\end{equation}

Now, we can apply the chain rule to $f=\overline{\log}$ (since $\left[x^{0}\right]g=0$), and
thus obtain
\begin{equation}
\left(\overline{\log}\circ g\right)^{\prime}=\left(\underbrace{\overline
{\log}^{\prime}}_{=\left(1+x\right)^{-1}}\circ\,g\right)\cdot g^{\prime
}=\left(\left(1+x\right)^{-1}\circ g\right)\cdot g^{\prime}.
\label{pf.prop.fps.exp-log-der.b.3}
\end{equation}

However, we claim that $\left(1+x\right)^{-1}\circ g=\left(1+g\right)
^{-1}$. Indeed, $\dfrac{1}{1+x}\circ g=\dfrac{1\circ
g}{\left(1+x\right)\circ g}$ (since the FPS $1+x$ is invertible), and
$1\circ\,g=\underline{1}$ and
$\left(1+x\right)\circ g=1+g$,
so $\left(1+x\right)^{-1}\circ g=\left(1+g\right)^{-1}$. Hence, (\ref{pf.prop.fps.exp-log-der.b.3})
becomes
\[
\left(\overline{\log}\circ g\right)^{\prime}=\left(1+g\right)^{-1}\cdot
g^{\prime}.
\]
This proves part \textbf{(b)}.
\end{proof}

\subsection{Helpers for Theorem~\ref{thm.fps.exp-log-inv}}

\begin{lemma}[ODE uniqueness: $h' = (1+h) \cdot g$]
\label{lem.fps.ode-uniqueness-mul-inv}
\lean{PowerSeries.eq_of_derivative_eq_mul_of_inv}
\leanhelper
\leanok
Let $h_1, h_2, g \in K\llbracket x \rrbracket$.
If $h_1' = (1 + h_1) \cdot g$ and $h_2' = (1 + h_2) \cdot g$
and $[x^0] h_1 = [x^0] h_2$, then $h_1 = h_2$.
\end{lemma}

\begin{proof}
By strong induction on the coefficient index.
The base case uses the matching constant terms.
For the inductive step, the ODE relates $[x^{n+1}] h \cdot (n+1)$ to a
convolution involving only lower coefficients; by the induction hypothesis,
these coincide for $h_1$ and $h_2$, and one cancels $n+1$ using the
$\mathbb{Q}$-algebra structure.
\end{proof}

\begin{lemma}[ODE uniqueness: $h' = 1$]
\label{lem.fps.ode-uniqueness-one}
\lean{PowerSeries.eq_of_derivative_eq_one}
\leanhelper
\leanok
Let $h_1, h_2 \in K\llbracket x \rrbracket$.
If $h_1' = 1$ and $h_2' = 1$ and $[x^0] h_1 = [x^0] h_2$,
then $h_1 = h_2$.
\end{lemma}

\begin{proof}
As in Lemma~\ref{lem.fps.ode-uniqueness-mul-inv}, by strong induction on
the coefficient index, cancelling $n+1$ at each step.
\end{proof}

\begin{lemma}[$\iota \circ \overline{\exp}$ times $\exp$]
\label{lem.fps.invOnePlusX-subst-expbar-mul-exp}
\lean{PowerSeries.invOnePlusX_subst_expbar_mul_exp}
\leanhelper
\leanok
We have $(\iota \circ \overline{\exp}) \cdot \exp = 1$.
\end{lemma}

\begin{proof}
We substitute $\overline{\exp}$ into the identity
$\iota \cdot (1 + x) = 1$, noting that
$(1 + x) \circ \overline{\exp} = 1 + \overline{\exp} = \exp$.
\end{proof}

\begin{lemma}
\label{lem.fps.compos-cst-term-0}
\lean{PowerSeries.constantCoeff_subst_of_constantCoeff_zero}
\leantarget
\leanok
Let $f,g\in K\left[\left[x\right]
\right]$ be two FPSs with $\left[x^{0}\right]g=0$. Then, $\left[
x^{0}\right]\left(f\circ g\right)=\left[x^{0}\right]f$.
\end{lemma}

\begin{proof}
\leanok
Write $f$ in the form
$f=\sum_{n\in\mathbb{N}}f_{n}x^{n}$ with $f_{0},f_{1},f_{2},\ldots\in K$.
Thus, $f_{0}=\left[x^{0}\right]f$. Now, $f\left[g\right]=\sum_{n\in\mathbb{N}}f_{n}g^{n}$. However,
$\left[x^{0}\right]
\left(\sum_{n\in\mathbb{N}}f_{n}g^{n}\right)=f_{0}=\left[x^{0}\right]
f$. In view of $f\circ g=f\left[g\right]=\sum_{n\in\mathbb{N}}f_{n}g^{n}$,
we can rewrite this as $\left[x^{0}\right]\left(f\circ g\right)
=\left[x^{0}\right]f$.
\end{proof}

\begin{theorem}
\label{thm.fps.exp-log-inv}
\lean{PowerSeries.expbar_comp_logbar, PowerSeries.logbar_comp_expbar}
\leantarget
\leanok
We have
\[
\overline{\exp}\circ\overline{\log}=x\ \ \ \ \ \ \ \ \ \ \text{and}%
\ \ \ \ \ \ \ \ \ \ \overline{\log}\circ\overline{\exp}=x.
\]
\end{theorem}

\begin{proof}
Let us first prove that
$\overline{\log}\circ\overline{\exp}=x$.

The idea of this proof is to show that $\overline{\log}\circ
\overline{\exp}$ and $x$ are two FPSs with the same constant term (namely,
$0$) and with the same derivative. Once this is proved, they must be equal
(since a FPS is determined by its constant term and its derivative over a $\mathbb{Q}$-algebra).

We have $\left[x^{0}\right]\overline{\exp}=0$
(since $\overline{\exp}=\sum_{n\geq1}\dfrac{1}{n!}x^{n}$). Hence, Lemma
\ref{lem.fps.compos-cst-term-0} (applied to $f=\overline{\log}$ and
$g=\overline{\exp}$) yields $\left[x^{0}\right]\left(\overline{\log
}\circ\overline{\exp}\right)=\left[x^{0}\right]\overline{\log}=0$ (since
$\overline{\log}=\sum_{n\geq1}\dfrac{\left(-1\right)^{n-1}}{n}x^{n}$).
Now,
\begin{equation}
\left[x^{0}\right]\left(\overline{\log}\circ\overline{\exp}-x\right)
=\underbrace{\left[x^{0}\right]\left(\overline{\log}\circ\overline{\exp
}\right)}_{=0}-\underbrace{\left[x^{0}\right]x}_{=0}=0.
\label{pf.thm.fps.exp-log-inv.a.1}
\end{equation}
However, $\overline{\exp}=\exp-1$ and thus $1+\overline{\exp}=\exp$. Now,
Proposition \ref{prop.fps.exp-log-der} \textbf{(b)} (applied to $g=\overline
{\exp}$) yields
\[
\left(\overline{\log}\circ\overline{\exp}\right)^{\prime}=\left(
\underbrace{1+\overline{\exp}}_{=\exp}\right)^{-1}\cdot\underbrace{\overline
{\exp}^{\prime}}_{=\exp}=\exp^{-1}\cdot\exp=1=x^{\prime}
\]
(where the underbrace uses (\ref{pf.prop.fps.exp-log-der.a.3}), and $x^{\prime}=1$). Hence,
$\overline{\log}\circ\overline{\exp}-x$ is constant (since its derivative is $0$). In view of
(\ref{pf.thm.fps.exp-log-inv.a.1}), this constant must be $0$. Thus, $\overline{\log
}\circ\overline{\exp}=x$.

Now it remains to prove that $\overline{\exp}\circ\overline{\log}=x$.

We shall first show
that $\exp\circ\overline{\log}=1+x$.

To wit: The FPS $1+x$ is invertible. Applying the quotient rule
to $f=\exp\circ\,\overline{\log}$ and
$g=1+x$, we obtain
\[
\left(\dfrac{\exp\circ\,\overline{\log}}{1+x}\right)^{\prime}
=\dfrac{\left(\exp\circ\,\overline{\log}\right)^{\prime}\cdot\left(
1+x\right)-\left(\exp\circ\,\overline{\log}\right)\cdot\left(
1+x\right)^{\prime}}{\left(1+x\right)^{2}}.
\]
In view of
\[
\left(\exp\circ\,\overline{\log}\right)^{\prime}
=\left(\exp\circ\,\overline{\log}\right)\cdot\underbrace{\overline{\log}^{\prime}
}_{=\left(1+x\right)^{-1}}
=\left(\exp\circ\,\overline{\log}\right)\cdot\left(1+x\right)^{-1}
\]
(where we used (\ref{pf.prop.fps.exp-log-der.b.2}))
and $\left(1+x\right)^{\prime}=1$, we can rewrite this as
\begin{align*}
\left(\dfrac{\exp\circ\,\overline{\log}}{1+x}\right)^{\prime}
=\dfrac{\left(\exp\circ\,\overline{\log}\right)-\left(\exp
\circ\,\overline{\log}\right)}{\left(1+x\right)^{2}}=0=0^{\prime}.
\end{align*}
Thus,
$\dfrac{\exp\circ\,\overline{\log}}{1+x}$ is constant, say equal to $\underline{a}$ for some $a\in K$.
Then $\exp\circ\,\overline{\log}=a\left(1+x\right)$, so
$\left[x^{0}\right]\left(\exp\circ\,\overline{\log}\right)=a$.
However, it is easy to see that $\left[x^{0}\right]\left(\exp
\circ\,\overline{\log}\right)=1$
(using Lemma \ref{lem.fps.compos-cst-term-0}).
Comparing, we find $a=1$. Thus, $\exp\circ\,\overline{\log}=1+x$.

Now, $\overline{\exp}=\exp-1=\exp+\,\underline{-1}$. Hence,
\begin{align*}
\overline{\exp}\circ\overline{\log}
=\underbrace{\exp\circ\,\overline{\log}}_{=1+x}+\underbrace{\underline{-1}
\circ\overline{\log}}_{=\underline{-1}}
=1+x+\underline{-1}=x.
\end{align*}
This completes the proof.
\end{proof}

\subsection{The exponential and the logarithm of an FPS}

\begin{definition}
\label{def.fps.Exp-Log-maps}
\lean{PowerSeries.PowerSeries₀, PowerSeries.PowerSeries₁, PowerSeries.Exp, PowerSeries.Log}
\leantarget
\leanok
\textbf{(a)} We let $K\left[\left[x\right]
\right]_{0}$ denote the set of all FPSs $f\in K\left[\left[x\right]
\right]$ with $\left[x^{0}\right]f=0$. \medskip

\textbf{(b)} We let $K\left[\left[x\right]\right]_{1}$ denote the set
of all FPSs $f\in K\left[\left[x\right]\right]$ with $\left[
x^{0}\right]f=1$. \medskip

\textbf{(c)} We define two maps
\begin{align*}
\operatorname{Exp}:K\left[\left[x\right]\right]_{0} &  \rightarrow
K\left[\left[x\right]\right]_{1},\\
g &  \mapsto\exp\circ g
\end{align*}
and
\begin{align*}
\operatorname{Log}:K\left[\left[x\right]\right]_{1} &  \rightarrow
K\left[\left[x\right]\right]_{0},\\
f &  \mapsto\overline{\log}\circ\left(f-1\right).
\end{align*}
(These two maps are well-defined according to parts \textbf{(c)} and
\textbf{(d)} of Lemma \ref{lem.fps.Exp-Log-maps-wd} below.)
\end{definition}

\begin{lemma}
\label{lem.fps.Exp-Log-maps-wd}
\lean{PowerSeries.PowerSeries₀.subst_mem, PowerSeries.PowerSeries₁.subst_mem, PowerSeries.exp_subst_mem_PowerSeries₁, PowerSeries.logbar_subst_sub_one_mem_PowerSeries₀}
\leantarget
\leanok
\textbf{(a)} For any $f,g\in K\left[\left[
x\right]\right]_{0}$, we have $f\circ g\in K\left[\left[x\right]
\right]_{0}$. \medskip

\textbf{(b)} For any $f\in K\left[\left[x\right]\right]_{1}$ and $g\in
K\left[\left[x\right]\right]_{0}$, we have $f\circ g\in K\left[
\left[x\right]\right]_{1}$. \medskip

\textbf{(c)} For any $g\in K\left[\left[x\right]\right]_{0}$, we have
$\exp\circ g\in K\left[\left[x\right]\right]_{1}$. \medskip

\textbf{(d)} For any $f\in K\left[\left[x\right]\right]_{1}$, we have
$f-1\in K\left[\left[x\right]\right]_{0}$ and $\overline{\log}
\circ\left(f-1\right)\in K\left[\left[x\right]\right]_{0}$.
\end{lemma}

\begin{proof}
\textbf{(a)} Let $f,g\in
K\left[\left[x\right]\right]_{0}$. Then $\left[
x^{0}\right]f=0$ and $\left[x^{0}\right]g=0$. Hence, Lemma
\ref{lem.fps.compos-cst-term-0} yields $\left[x^{0}\right]\left(f\circ
g\right)=\left[x^{0}\right]f=0$. In other words, $f\circ g\in K\left[
\left[x\right]\right]_{0}$.

\textbf{(b)} This is analogous to the proof of part \textbf{(a)}.

\textbf{(c)} Let $g\in K\left[\left[x\right]\right]_{0}$. From
$\exp=\sum_{n\in\mathbb{N}}\dfrac{1}{n!}x^{n}$, we obtain $\left[
x^{0}\right]\exp=1$, so that $\exp\in K\left[\left[
x\right]\right]_{1}$. Hence, part
\textbf{(b)} (applied to $f=\exp$) yields $\exp\circ g\in K\left[\left[
x\right]\right]_{1}$.

\textbf{(d)} Let $f\in K\left[\left[x\right]\right]_{1}$. Thus,
$\left[x^{0}\right]f=1$. Now,
$\left[x^{0}\right]\left(f-1\right)=\left[x^{0}\right]
f-\left[x^{0}\right]1=1-1=0$, so that $f-1\in
K\left[\left[x\right]\right]_{0}$. Furthermore, $\left[x^{0}\right]
\overline{\log}=0$ and thus $\overline{\log}\in K\left[\left[
x\right]\right]_{0}$. Hence, part
\textbf{(a)} (applied to $\overline{\log}$ and $f-1$ instead of $f$ and $g$)
yields $\overline{\log}\circ\left(f-1\right)\in K\left[\left[x\right]
\right]_{0}$.
\end{proof}

\begin{lemma}
\label{lem.fps.Exp-Log-maps-inv}
\lean{PowerSeries.Log_Exp, PowerSeries.Exp_Log, PowerSeries.Exp_Log_inverse}
\leantarget
\leanok
The maps $\operatorname{Exp}$ and
$\operatorname{Log}$ are mutually inverse bijections between $K\left[
\left[x\right]\right]_{0}$ and $K\left[\left[x\right]\right]_{1}$.
\end{lemma}

\begin{proof}
\leanok
First, we make a simple
auxiliary observation: Each $g\in K\left[\left[x\right]\right]_{0}$
satisfies
\begin{equation}
\exp\circ g=\overline{\exp}\circ g+1. \label{pf.lem.fps.Exp-Log-maps-inv.1}
\end{equation}

[\textit{Proof of (\ref{pf.lem.fps.Exp-Log-maps-inv.1}):} Recall that $\overline{\exp}=\exp-1$, so
that $\exp=\overline{\exp}+\underline{1}$. Hence,
$\exp\circ g=\left(\overline{\exp}+\underline{1}\right)\circ g=\overline
{\exp}\circ g+\underline{1}\circ g$. But
$\underline{1}\circ g=\underline{1}=1$. Hence, $\exp\circ g=\overline{\exp
}\circ g+1$.]

Now, let us show that $\operatorname{Exp}\circ\operatorname{Log}
=\operatorname{id}$. Indeed, we fix some $f\in K\left[\left[x\right]
\right]_{1}$. Then, $f-1\in K\left[\left[x\right]\right]_{0}$ (by
Lemma \ref{lem.fps.Exp-Log-maps-wd} \textbf{(d)}). By associativity of composition and Theorem \ref{thm.fps.exp-log-inv}, we get
\[
\overline{\exp}\circ\left(\overline{\log}\circ\left(f-1\right)\right)
=\underbrace{\left(\overline{\exp}\circ\overline{\log}\right)
}_{=x}\circ\left(f-1\right)=x\circ\left(f-1\right)
=f-1.
\]
Hence
\begin{align*}
\left(\operatorname{Exp}\circ\operatorname{Log}\right)\left(f\right)
&  =\exp\circ\left(\operatorname{Log}f\right) \\
&  =\overline{\exp}\circ\underbrace{\left(\operatorname{Log}f\right)
}_{=\overline{\log}\circ\left(f-1\right)}+\,1
\end{align*}
(by (\ref{pf.lem.fps.Exp-Log-maps-inv.1})), so
\begin{align*}
\left(\operatorname{Exp}\circ\operatorname{Log}\right)\left(f\right)
&  =\underbrace{\overline{\exp}\circ\left(\overline{\log}\circ\left(
f-1\right)\right)}_{=f-1}+\,1
=\left(f-1\right)+1=f.
\end{align*}

Using a similar argument (with $\overline{\log}\circ\overline{\exp}=x$ from Theorem \ref{thm.fps.exp-log-inv}), we can show that $\operatorname{Log}\circ
\operatorname{Exp}=\operatorname{id}$.
Combining, $\operatorname{Exp}$ and
$\operatorname{Log}$ are mutually inverse bijections.
\end{proof}

\subsection{Helpers for the group structure}

\begin{lemma}[$f - 1 \in K\llbracket x \rrbracket_0$ when $f \in K\llbracket x \rrbracket_1$]
\label{lem.fps.sub-one-mem-PS0}
\lean{PowerSeries.sub_one_mem_PowerSeries₀}
\leanhelper
\leanok
If $f \in K\llbracket x \rrbracket_1$ (i.e., $[x^0]f = 1$),
then $f - 1 \in K\llbracket x \rrbracket_0$ (i.e., $[x^0](f-1) = 0$).
\end{lemma}

\begin{proof}
\leanok
$[x^0](f - 1) = [x^0]f - 1 = 1 - 1 = 0$.
\end{proof}

\begin{lemma}[ODE uniqueness: $h' = h \cdot g$]
\label{lem.fps.ode-uniqueness-mul-self}
\lean{PowerSeries.eq_of_derivative_eq_mul_self}
\leanhelper
\leanok
Let $h_1, h_2, g \in K\llbracket x \rrbracket$.
If $h_1' = h_1 \cdot g$ and $h_2' = h_2 \cdot g$
and $[x^0] h_1 = [x^0] h_2$, then $h_1 = h_2$.
\end{lemma}

\begin{proof}
By strong induction on the coefficient index. The ODE gives
$[x^{n+1}] h \cdot (n+1) = [x^n](h \cdot g)$; the right side is a
convolution involving only coefficients of $h$ of index $\le n$,
which agree by induction hypothesis. Cancelling $n+1$ uses the
$\mathbb{Q}$-algebra structure.
\end{proof}

\begin{lemma}[$\operatorname{Exp}(0) = 1$ and $\operatorname{Log}(1) = 0$]
\label{lem.fps.Exp-zero-Log-one}
\lean{PowerSeries.Exp_zero, PowerSeries.Log_one}
\leanhelper
\leanok
The map $\operatorname{Exp}$ sends $0 \in K\llbracket x \rrbracket_0$
to $1 \in K\llbracket x \rrbracket_1$, and $\operatorname{Log}$
sends $1$ to $0$.
\end{lemma}

\begin{proof}
\leanok
$\operatorname{Exp}(0) = \exp \circ 0 = 1$ (since $[x^0]\exp = 1$ and
substituting $0$ picks out the constant term).
$\operatorname{Log}(1) = \overline{\log} \circ (1 - 1) = \overline{\log} \circ 0 = 0$
(since $[x^0]\overline{\log} = 0$).
\end{proof}

\subsection{Addition to multiplication}

\begin{lemma}
\label{lem.fps.Exp-Log-additive}
\lean{PowerSeries.Exp_add, PowerSeries.Log_mul}
\leantarget
\leanok
\textbf{(a)} For any $f,g\in K\left[\left[
x\right]\right]_{0}$, we have
\[
\operatorname{Exp}\left(f+g\right)=\operatorname{Exp}f\cdot
\operatorname{Exp}g.
\]

\textbf{(b)} For any $f,g\in K\left[\left[x\right]\right]_{1}$, we
have
\[
\operatorname{Log}\left(fg\right)=\operatorname{Log}
f+\operatorname{Log}g.
\]
\end{lemma}

\begin{proof}
\textbf{(a)} Let $f,g\in K\left[\left[x\right]\right]_{0}$. Thus, $\left[
x^{0}\right]f=0$ and $\left[x^{0}\right]g=0$. By the definition of $\operatorname{Exp}$, we have
\[
\operatorname{Exp}f=\exp\circ f=\sum_{n\in\mathbb{N}
}\dfrac{1}{n!}f^{n}.
\]
Similarly, $\operatorname{Exp}g=\sum_{n\in\mathbb{N}}\dfrac{1}{n!}g^{n}$
and
$\operatorname{Exp}\left(f+g\right)=\sum_{n\in\mathbb{N}}\dfrac{1}
{n!}\left(f+g\right)^{n}$.

Now, the latter equality becomes
\begin{align*}
\operatorname{Exp}\left(f+g\right) &  =\sum_{n\in\mathbb{N}}\dfrac{1}
{n!}\sum\limits_{k=0}
^{n}\dbinom{n}{k}f^{k}g^{n-k}
=\sum_{k\in\mathbb{N}}\ \ \sum\limits_{\ell
\in\mathbb{N}}\dfrac{1}{k!\ell!}f^{k}g^{\ell}.
\end{align*}
Comparing this with
$\operatorname{Exp}f\cdot\operatorname{Exp}g
=\sum_{k\in\mathbb{N}}\ \ \sum_{\ell\in
\mathbb{N}}\dfrac{1}{k!\ell!}f^{k}g^{\ell}$,
we obtain $\operatorname{Exp}\left(f+g\right)=\operatorname{Exp}
f\cdot\operatorname{Exp}g$.

(In the Lean formalization, this is proved via ODE uniqueness:
both sides satisfy the ODE $h^{\prime} = h \cdot (f^{\prime} + g^{\prime})$ with the same initial condition.)

\textbf{(b)} This follows from part \textbf{(a)}, since we know that
$\operatorname{Log}$ is inverse to $\operatorname{Exp}$. Setting
$u=\operatorname{Log}f$ and $v=\operatorname{Log}g$, part \textbf{(a)}
gives $\operatorname{Exp}\left(u+v\right)=\operatorname{Exp}u\cdot\operatorname{Exp}v$,
and applying $\operatorname{Log}$ (using Lemma \ref{lem.fps.Exp-Log-maps-inv}) yields
$\operatorname{Log}\left(fg\right)=\operatorname{Log}f+\operatorname{Log}g$.
\end{proof}

\begin{proposition}
\label{prop.fps.Exp-Log-groups}
\lean{PowerSeries.PowerSeries₀.addSubgroup, PowerSeries.PowerSeries₁.subgroup}
\leantarget
\leanok
\textbf{(a)} The subset $K\left[\left[
x\right]\right]_{0}$ of $K\left[\left[x\right]\right]$ is closed
under addition and subtraction and contains $0$, and thus forms a group
$\left(K\left[\left[x\right]\right]_{0},+,0\right)$. \medskip

\textbf{(b)} The subset $K\left[\left[x\right]\right]_{1}$ of
$K\left[\left[x\right]\right]$ is closed under multiplication and
division and contains $1$, and thus forms a group $\left(K\left[\left[
x\right]\right]_{1},\cdot,1\right)$.
\end{proposition}

\begin{proof}
\textbf{(a)} It is clear
that the set $K\left[\left[x\right]\right]_{0}$ contains the FPS $0$
(since $\left[x^{0}\right]0=0$). If $f,g\in K\left[\left[x\right]\right]_{0}$, then
$\left[x^{0}\right]f=0$ and $\left[x^{0}\right]g=0$, and therefore
$f+g\in K\left[\left[x\right]\right]_{0}$ (since
$\left[x^{0}\right]\left(
f+g\right)=0+0=0$) and $f-g\in K\left[\left[x\right]\right]
_{0}$ (by a similar argument).

\textbf{(b)} Any $a\in K\left[\left[
x\right]\right]_{1}$ is invertible
in $K\left[\left[x\right]\right]$ (since $\left[x^{0}\right]a=1$ is invertible in $K$).
Next, $K\left[\left[x\right]\right]_{1}$ is closed under
multiplication: if $f,g\in K\left[\left[x\right]\right]_{1}$,
then $\left[x^{0}\right]f=1$ and $\left[x^{0}\right]g=1$, and
therefore $\left[x^{0}\right]\left(
fg\right)=1\cdot 1=1$.
Similarly, $K\left[\left[x\right]\right]_{1}$ is closed
under division.
\end{proof}

\begin{theorem}
\label{thm.fps.Exp-Log-group-iso}
\lean{PowerSeries.Exp_Log_groupIso}
\leantarget
\leanok
The maps
\[
\operatorname{Exp}:\left(K\left[\left[x\right]\right]_{0}
,+,0\right)\rightarrow\left(K\left[\left[x\right]\right]_{1}
,\cdot,1\right)
\]
and
\[
\operatorname{Log}:\left(K\left[\left[x\right]\right]_{1}
,\cdot,1\right)\rightarrow\left(K\left[\left[x\right]\right]
_{0},+,0\right)
\]
are mutually inverse group isomorphisms.
\end{theorem}

\begin{proof}
Lemma
\ref{lem.fps.Exp-Log-additive} yields that these two maps are group
homomorphisms. Lemma \ref{lem.fps.Exp-Log-maps-inv}
shows that they are mutually inverse. Combining these results, we conclude
that these two maps are mutually inverse group isomorphisms.
\end{proof}

\subsection{Helpers for the logarithmic derivative}

\begin{lemma}[$\iota = (1+x)^{-1}$]
\label{lem.fps.invOnePlusX-eq-inv}
\lean{PowerSeries.invOnePlusX_eq_inv}
\leanhelper
\leanok
Over a field (or $\mathbb{Q}$-algebra), $\iota$
equals the formal inverse $(1 + x)^{-1}$ in $K\llbracket x \rrbracket$.
\end{lemma}

\begin{proof}
\leanok
Both sides multiply with $(1+x)$ to give $1$; since the constant
coefficient of $1+x$ is nonzero, the inverse is unique.
\end{proof}

\begin{lemma}[Substitution of inverses]
\label{lem.fps.subst-inv}
\lean{PowerSeries.subst_inv_of_mul_eq_one}
\leanhelper
\leanok
If $a \cdot b = 1$ in $K\llbracket x \rrbracket$
and $g$ has constant term $0$, then $(a \circ g)^{-1} = b \circ g$.
\end{lemma}

\begin{proof}
\leanok
Substitution is a ring homomorphism, so $(a \circ g) \cdot (b \circ g) = (a \cdot b) \circ g = 1 \circ g = 1$.
Since $[x^0](a \circ g) \ne 0$, the inverse is unique.
\end{proof}

\subsection{The logarithmic derivative}

\begin{definition}
\label{def.fps.loder.1}
\lean{PowerSeries.loder}
\leantarget
\leanok
In this definition, we do \textbf{not} use Convention
\ref{conv.fps.exp.K-Q-alg}; thus, $K$ can be an arbitrary commutative ring.
However, we set $K\left[\left[x\right]\right]_{1}=\left\{f\in
K\left[\left[x\right]\right]\ \mid\ \left[x^{0}\right]f=1\right\}
$.

For any FPS $f\in K\left[\left[x\right]\right]_{1}$, we define the
\emph{logarithmic derivative} $\operatorname{loder}f\in K\left[\left[
x\right]\right]$ to be the FPS $\dfrac{f^{\prime}}{f}$. (This is
well-defined, since $f$ is easily seen to be invertible.)
\end{definition}

\begin{lemma}[Invertibility of FPS with constant term $1$]
\label{lem.fps.isUnit-constantCoeff-one}
\lean{PowerSeries.isUnit_of_constantCoeff_eq_one, PowerSeries.isUnit_iff_constantCoeff_isUnit, PowerSeries.mul_inv_cancel_of_constantCoeff_eq_one, PowerSeries.inv_mul_cancel_of_constantCoeff_eq_one, PowerSeries.constantCoeff_inv_of_constantCoeff_eq_one}
\leanhelper
\leanok
If $[x^0]f = 1$, then $f$ is a unit in $K\llbracket x \rrbracket$, i.e.,
$f \cdot f^{-1} = 1$ and $f^{-1} \cdot f = 1$.
Moreover, $[x^0](f^{-1}) = 1$.
More generally, $f$ is a unit iff $[x^0]f$ is a unit in $K$.
\end{lemma}

\begin{proof}
\leanok
Since $[x^0]f$ is a unit, the invertibility criterion for power series
gives the result (a power series is invertible if and only if its constant coefficient is invertible).
The constant coefficient of $f^{-1}$ is the inverse of $[x^0]f = 1$.
\end{proof}

\begin{lemma}[$\operatorname{loder}(1) = 0$]
\label{lem.fps.loder-one}
\lean{PowerSeries.loder_one}
\leanhelper
\leanok
$\operatorname{loder}(1) = 0$.
\end{lemma}

\begin{proof}
\leanok
$1' = 0$, so $\operatorname{loder}(1) = 0 \cdot 1^{-1} = 0$.
\end{proof}

\begin{proposition}
\label{prop.fps.loder.log}
\lean{PowerSeries.loder_eq_derivative_Log}
\leantarget
\leanok
Let $K$ be a commutative $\mathbb{Q}$-algebra. Let
$f\in K\left[\left[x\right]\right]_{1}$ be an FPS. Then,
$\operatorname{loder}f=\left(\operatorname{Log}f\right)^{\prime}$.
\end{proposition}

\begin{proof}
\leanok
The definition of
$\operatorname{Log}$ yields $\operatorname{Log}f=\overline{\log}\circ\left(
f-1\right)$.

From $f\in K\left[\left[x\right]\right]_{1}$, we obtain $\left[
x^{0}\right]f=1$. Now, $\left[x^{0}\right]
\left(f-1\right)=\left[x^{0}\right]f-\left[x^{0}\right]1=0$. Hence, Proposition
\ref{prop.fps.exp-log-der} \textbf{(b)} (applied to $g=f-1$) yields
\[
\left(\overline{\log}\circ\left(f-1\right)\right)^{\prime}=\left(
\underbrace{1+\left(f-1\right)}_{=f}\right)^{-1}\cdot\left(f-1\right)
^{\prime}=f^{-1}\cdot\left(f-1\right)^{\prime}=\dfrac{\left(f-1\right)
^{\prime}}{f}.
\]
In view of $\operatorname{Log}f=\overline{\log}\circ\left(f-1\right)$, we
can rewrite this as
$\left(\operatorname{Log}f\right)^{\prime}=\dfrac{\left(f-1\right)
^{\prime}}{f}$.

Since $\left(f-1\right)^{\prime}=f^{\prime}$, we can rewrite
this further as $\left(\operatorname{Log}f\right)^{\prime}=\dfrac
{f^{\prime}}{f}=\operatorname{loder}f$.
\end{proof}

\begin{proposition}
\label{prop.fps.loder.prod}
\lean{PowerSeries.loder_mul}
\leantarget
\leanok
Let $f,g\in K\left[\left[x\right]\right]
_{1}$ be two FPSs. Then, $\operatorname{loder}\left(fg\right)
=\operatorname{loder}f+\operatorname{loder}g$.

(Here, we do \textbf{not} use Convention \ref{conv.fps.exp.K-Q-alg}; thus, $K$
can be an arbitrary commutative ring.)
\end{proposition}

\begin{proof}
The definition of a
logarithmic derivative yields $\operatorname{loder}f=\dfrac{f^{\prime}}{f}$
and $\operatorname{loder}g=\dfrac{g^{\prime}}{g}$, but it also yields
\begin{align*}
\operatorname{loder}\left(fg\right) &  =\dfrac{\left(fg\right)
^{\prime}}{fg}=\dfrac{f^{\prime}g+fg^{\prime}}{fg}
\qquad\text{(by the product rule)}
=\underbrace{\dfrac{f^{\prime}}{f}}_{=\operatorname{loder}f}
+\underbrace{\dfrac{g^{\prime}}{g}}_{=\operatorname{loder}g}
=\operatorname{loder}f+\operatorname{loder}g.
\end{align*}
\end{proof}

\begin{corollary}
\label{cor.fps.loder.prodk}
\lean{PowerSeries.loder_prod}
\leantarget
\leanok
Let $f_{1},f_{2},\ldots,f_{k}$ be any $k$ FPSs in
$K\left[\left[x\right]\right]_{1}$. Then,
\[
\operatorname{loder}\left(f_{1}f_{2}\cdots f_{k}\right)
=\operatorname{loder}\left(f_{1}\right)+\operatorname{loder}\left(
f_{2}\right)+\cdots+\operatorname{loder}\left(f_{k}\right).
\]

(Here, we do \textbf{not} use Convention \ref{conv.fps.exp.K-Q-alg}; thus, $K$
can be an arbitrary commutative ring.)
\end{corollary}

\begin{proof}
Induct on $k$. The \textit{base
case} ($k=0$) requires showing that $\operatorname{loder}1=0$, which is easy
to see from the definition of a logarithmic derivative (since $1^{\prime}=0$).
The \textit{induction step} follows easily from Proposition
\ref{prop.fps.loder.prod}.
\end{proof}

\begin{corollary}
\label{cor.fps.loder.inv}
\lean{PowerSeries.loder_inv}
\leantarget
\leanok
Let $f$ be any FPS in $K\left[\left[x\right]
\right]_{1}$. Then, $\operatorname{loder}\left(f^{-1}\right)
=-\operatorname{loder}f$.

(Here, we do \textbf{not} use Convention \ref{conv.fps.exp.K-Q-alg}; thus, $K$
can be an arbitrary commutative ring.)
\end{corollary}

\begin{proof}
First of all, $f$ is invertible. This shows that $f^{-1}$
is well-defined.

Proposition \ref{prop.fps.loder.prod} (applied to $g=f^{-1}$) yields
$\operatorname{loder}\left(ff^{-1}\right)=\operatorname{loder}
f+\operatorname{loder}\left(f^{-1}\right)$. However, we also have
$\operatorname{loder}\left(\underbrace{ff^{-1}}_{=1}\right)
=\operatorname{loder}1=0$ (since $1^{\prime}=0$). Comparing these two
equalities, we obtain $\operatorname{loder}f+\operatorname{loder}\left(
f^{-1}\right)=0$. In other words, $\operatorname{loder}\left(
f^{-1}\right)=-\operatorname{loder}f$.
\end{proof}

\subsection{Summable and multipliable families}

\begin{definition}[Summable and multipliable families]
\label{def.fps.summable-multipliable}
\lean{PowerSeries.SummableFPS₀, PowerSeries.MultipliableFPS₁, PowerSeries.summableFPS₀Sum, PowerSeries.multipliableFPS₁Prod}
\leanhelper
\leanok
A family $(f_i)_{i \in I}$ of FPSs in $K\llbracket x \rrbracket_0$
is \emph{summable} if for each coefficient index $n$, only finitely many
$[x^n](f_i)$ are nonzero.

A family $(f_i)_{i \in I}$ of FPSs in $K\llbracket x \rrbracket_1$
is \emph{multipliable} if for each coefficient index $n$, all but finitely many
$f_i$ satisfy $[x^k](f_i - 1) = 0$ for all $1 \le k \le n$.

For summable families, the coefficient-wise sum $\sum_{i \in I} f_i$
belongs to $K\llbracket x \rrbracket_0$.
For multipliable families, the product $\prod_{i \in I} f_i$
belongs to $K\llbracket x \rrbracket_1$.
\end{definition}

\begin{proof}
For the sum: the constant coefficient of each $f_i$ is $0$, so the sum
of constant coefficients is $0$.
For the product: the constant coefficient of $\prod f_i$ is $\prod [x^0](f_i) = \prod 1 = 1$.
\end{proof}

\subsection{Exp and Log for infinite products and sums}

\begin{lemma}[Summability of $\operatorname{Log}$-images]
\label{lem.fps.Log-summable}
\lean{PowerSeries.Log_summable_of_multipliable}
\leanhelper
\leanok
If $(f_i)_{i \in I}$ is a multipliable family in
$K\llbracket x \rrbracket_1$, then $(\operatorname{Log} f_i)_{i \in I}$
is a summable family in $K\llbracket x \rrbracket_0$
(meaning that for each coefficient index $n$, only finitely many
$[\operatorname{Log} f_i]_n$ are nonzero).
\end{lemma}

\begin{proof}
For each $n$, let $M$ be an $x^n$-approximator for the product family.
If $i \notin M$, then $[x^k](f_i) = 0$ for $1 \le k \le n$
(by the approximator property).
This forces $[(f_i - 1)]_k = 0$ for $0 \le k \le n$,
and a substitution argument shows $[\operatorname{Log} f_i]_n = 0$.
Hence the set of $i$ with nonzero $n$-th coefficient is contained in $M$.
\end{proof}

\begin{proposition}[$\operatorname{Log}$ of infinite product]
\label{prop.fps.Exp-Log-infprod}
\lean{PowerSeries.Log_tprod}
\leanhelper
\leanok
If $(f_i)_{i \in I}$ is a multipliable family in
$K\llbracket x \rrbracket_1$, then
\[
\operatorname{Log}\Bigl(\prod_{i \in I} f_i\Bigr)
= \sum_{i \in I} \operatorname{Log}(f_i),
\]
where the right side is the coefficient-wise sum
(which is well-defined by Lemma~\ref{lem.fps.Log-summable}).
\end{proposition}

\begin{proof}
\leanok
For each $n$, choose an $x^n$-approximator $M$.
By the finite version ($\operatorname{Log}$ of a finite product
equals the sum of $\operatorname{Log}$s, which follows from
Lemma~\ref{lem.fps.Exp-Log-additive}(b) by induction), we have
$\operatorname{Log}(\prod_{i \in M} f_i) = \sum_{i \in M} \operatorname{Log}(f_i)$.
The $n$-th coefficient of the infinite product agrees with that of
the finite product over $M$, and the $n$-th coefficient of the sum
is the finite sum over $M$ (since $[\operatorname{Log} f_i]_n = 0$
for $i \notin M$). A coefficient comparison argument completes the proof.
\end{proof}

\begin{lemma}[Multipliability of $\operatorname{Exp}$-images]
\label{lem.fps.Exp-multipliable}
\lean{PowerSeries.Exp_multipliable_of_summable}
\leanhelper
\leanok
If $(g_i)_{i \in I}$ is a summable family in
$K\llbracket x \rrbracket_0$, then $(\operatorname{Exp} g_i)_{i \in I}$
is a multipliable family in $K\llbracket x \rrbracket_1$.
\end{lemma}

\begin{proof}
\leanok
Since $\operatorname{Exp}(g_i) = 1 + \overline{\exp} \circ g_i$,
the family is of the form $(1 + h_i)$ where $h_i = \overline{\exp} \circ g_i$.
The key observation is that if $[x^k](g_i) = 0$ for all $k \le n$,
then $[x^n](\overline{\exp} \circ g_i) = 0$ as well
(since $\overline{\exp}$ has constant term $0$, so the substitution
preserves vanishing of low-order terms).
The summability of $(g_i)$ then implies the multipliability of
$(1 + h_i)$.
\end{proof}

\begin{proposition}[$\operatorname{Exp}$ of infinite sum]
\label{prop.fps.Exp-Log-infsum}
\lean{PowerSeries.Exp_sum}
\leanhelper
\leanok
If $(g_i)_{i \in I}$ is a summable family in
$K\llbracket x \rrbracket_0$, then
\[
\operatorname{Exp}\Bigl(\sum_{i \in I} g_i\Bigr)
= \prod_{i \in I} \operatorname{Exp}(g_i).
\]
\end{proposition}

\begin{proof}
\leanok
Let $f_i = \operatorname{Exp}(g_i)$.
By Proposition~\ref{prop.fps.Exp-Log-infprod},
$\operatorname{Log}(\prod f_i) = \sum \operatorname{Log}(f_i) = \sum g_i$
(using $\operatorname{Log} \circ \operatorname{Exp} = \mathrm{id}$).
Applying $\operatorname{Exp}$ to both sides and using
$\operatorname{Exp} \circ \operatorname{Log} = \mathrm{id}$ yields the result.
\end{proof}
